\begin{frame}{``분산이 크다''의 구체적 그림}
  진짜 경계가 대각선으로 갈라진 두 덩어리, 경계 근처에 \textbf{한
  점만} 잘못 찍혀 있다고 하자:
  \vskip0.8em
  \begin{columns}[T]
    \begin{column}{0.48\textwidth}
      \kb{$k=1$}
      \begin{itemize}\small
        \item 잘못 찍힌 점 바로 옆에서 ``그 점이 정답''
        \item \textbf{경계 전체가 한 점의 위치를 따라 출렁}
      \end{itemize}
    \end{column}
    \begin{column}{0.48\textwidth}
      \kb{$k=5$}
      \begin{itemize}\small
        \item 잘못 찍힌 한 점을 4개 정상점이 \textbf{투표로 이김}
        \item 그 한 점은 경계에 거의 영향 없음
      \end{itemize}
    \end{column}
  \end{columns}
  \vskip0.8em
  반대로 진짜 경계가 \textbf{복잡하게 꼬여} 있으면, $k=5$는
  ``이웃 5개 다수결''로 세부 곡선을 따라가지 못하고 매끄럽게
  그은다 --- 이쪽은 \kb{편향}의 문제.
  \vskip0.6em
  \begin{block}{정의의 kNN 번역}
    분산이 크다 = \textbf{훈련 데이터의 미세한 변화}(노이즈 한 점,
    표본 재추출)가 \textbf{예측의 미세한 변화}(경계의 출렁임)를
    낳는 정도. kNN에서는 이 인과가 \textbf{눈으로 직접 확인}
    가능 --- 신경망에서 ``분산''을 눈으로 보기 훨씬 어렵다.
  \end{block}
\end{frame}
